\subsection{零交換子CYの七つの必要対称射影型}
\label{subsec:cubic-CY-CI22-projection-atlas}
以下の $31,22$ 等は対称射影の基点スキームの型であり、中山作用のJordan型ではない。
\begin{theorem}\label{thm:cubic-CY-CI22-seven-projections}
定理\ref{thm:cubic-CY-zero-commutator-CI22}の対称射影は、線型生成元変換を除き、
次の相異なる七型で必要被覆される。
\[
\begin{array}{c|l|c}
\text{型}&U\text{ の基底}&\text{基点スキーム}\\\hline
1111&ab,ac,ad,bc,bd,cd&1+1+1+1\\
211&b^2,ac,ad,bc,bd,cd&2+1+1\\
22&b^2,d^2,ac,ad,bc,bd&2+2\\
31&ac-b^2,bc,c^2,ad,bd,cd&3+1\\
4&ac-b^2,ad-bc,bd,c^2,cd,d^2&4\\
L&a^2,ab,b^2,bc,ac-bd,ad&\text{被約直線}\\
D&a^2,ab,b^2,c^2,cd,d^2&\varnothing
\end{array}
\]
これは射影空間内の六二次式系の必要分類である。
その上のAS4係数・中央三次切断・H3許容性・全同型を同時に分類したものではない。
\end{theorem}
\begin{proof}
固定出典\cite{AS4CYpinned}の二次鉛筆射影定理の必要部分を用いる。
ここで使う幾何学的全分岐と、独立に再生した全有限恒等式は以下の通りである。

二次関係の鉛筆を $\langle Q,P\rangle$ とし、可逆な $Q$ の対称行列を $A$、
$P$ の行列を $H$、$T=A^{-1}H$ とする。
$T$ は $A$ に関して自己共役である。
像の一般滑らかな点 $q$ で、$Q$ の接二次空間 $q^\perp/kq$ 内の法線は $Tq$ の類である。
従って三次元接形式の退化条件は $q^tAT^2q=0$。
像が整な $CI(2,2)$ であり、その二次idealが鉛筆全体なので、恒等的に退化することは
\[
 T^2\in\langle I,T\rangle
\]
と同値である。$T$ はスカラーでないから、この場合の最小多項式は次数2となる。

二根が相異なれば、自己共役性により二つの固有空間は非退化で直交する。
その次元を $m,6-m$ とすると鉛筆にはrank $m,6-m$ の関係がある。
独立な二次式間の非零関係はrank1,2になれない（一次結合の積への因数分解が独立性に反する）ので、
$m=3$ である。
三式のrank3関係は $q_0q_2=q_1^2$ とでき、一意分解による原始vector平方の復元から
二つの一次式の全二次式系となる。
二群の一次式空間が交われば平方を共有し $\dim U<6$ となるので、両群は直和である。
従って $D$ を得る。

重根ならスカラーを引いて $N^2=0$、$N\ne0$ となる。
$AN$ は非零の二次関係なのでrankが少なくとも3、$N^2=0$ から高々3、従って正確に3である。
自己共役の双対基底により、鉛筆は非退化双曲形式とrank3関係へ正規化される。
rank3関係の三式を $a^2,ab,b^2$ とすれば、残る三二次式は
\[
 f_0a^2+f_1ab+f_2b^2=0
\]
を満たす。
その全二次syzygyは $(-b\ell,a\ell-bm,am)$、$\ell,m$ は一次式である。
係数写像はrank22、parameter写像はrank8であり、全解を尽くす。
$\ell,m$ の $\langle a,b\rangle$ を法とした類が従属なら $\dim U\le5$ となるので、
$(a,b,\ell,m)$ は座標を成す。これで $L$ を得る。
この系の基点は局所的に $(a,b)$ で定まり、被約直線である。
鉛筆は
\[
 Q=z_0z_3-z_1z_4-z_2z_5,\qquad P=z_0z_2-z_1^2
\]
で、$c\mapsto c+\alpha a+\beta b$, $d\mapsto d+\gamma a+\delta b$ は
$Q\mapsto Q+(\beta-\gamma)P$ と作用する。
従ってrank6関係の印付けも全て同じ軌道に入る。

常に退化する場合以外では非退化接形式の点があり、前節の対称行列 $M$ の表示を使える。
$b=(b_1,b_2,b_3)$、$B=b^tb$、$r=b^tMb$、$v=b\times Mb$ とする。
基点では $b\ne0$ であり、$b_i\ne0$ の開集合で $a=-v_i/b_i$ を消去すると、
残る全式は $B=r=0$ に一致する。
これは一般対称 $M$ の全係数のまま、次の三恒等式で確認できる：
\[
\begin{aligned}
 b_1v_2-b_2v_1&=b_3r-(Mb)_3B,\\
 b_1v_3-b_3v_1&=-b_2r+(Mb)_2B,\\
 v_1^2+b_1^2(Mb)^t(Mb)&=B((Mb)_2^2+(Mb)_3^2)+r(2b_1(Mb)_1-r).
\end{aligned}
\]
$M$ が非スカラーなので、基点は滑らかな円錐曲線 $B=0$ 上の長さ4のCartier切断であり、曲線的である。

この有限スキーム $Z\subset\mathbb P^3$ は超平面に含まれない。
実際、$\alpha a+\beta^tb$ が $Z$ 上で零なら、$\alpha=0$ は円錐曲線の長さ4切断を直線が含むことになり不可能。
$\alpha\ne0$ では $a=-\beta^tb$ とし、
$v-(\beta^tb)b=cB+er$ と書ける。
$b$ と内積を取り、互いに素な $B,r$ の次数を比較すると $e=0,c=-\beta$ となる。
従って $b\times Mb=b\times(b\times\beta)$ となり、交叉積の核がスカラー行列だけであることから、
対称部分 $M$ がスカラーとなって矛盾する。

従って $H^0(\mathbb P^3,\mathcal O(1))\to H^0(Z,\mathcal O_Z(1))$ は四次元同士の同型である。
$Z$ 上可逆な一次切断を一つ選べば、その乗算により次数二の制限も全射となる。
よって $\dim I_Z(2)=10-4=6$、$U=I_Z(2)$ である。
曲線的な長さ4の局所成分は $k[t]/(t^n)$ なので、分割は
$1111,211,22,31,4$ の五つだけである。
有限スキーム上の可逆層は自明であり、全切断による埋込みなので、
その抽象スキーム型が射影同値を決定する。
各直積有限環の基底 $1,t,\ldots$ に対する次数二の乗算kernelを取ると、表の最初の五行が得られる。
それぞれの全kernelが六次元であることを独立計算した。
五つの基点型、被約直線、空集合は区別されるので、七型は相異なる。
これで全場合を尽くす。
\end{proof}

独立再現器は\\
\texttt{verify\_cubic\_CY\_CI22\_projection\_atlas.py}。
七つの全積kernelと鉛筆行列式、$I,T,T^2$ の全rank、五有限環の乗算kernel、
全syzygyと $L$ のmarked関係変換、一般 $M$ の全基点恒等式、五根型の対称行列例を検証する。
これらは必要な対称射影の全被覆を支えるもので、各射影上の代数係数分類やAS十分性の検証ではない。
